% Generated draft: requires mathematical and source review.
\documentclass[11pt]{article}
\usepackage[margin=25mm]{geometry}
\usepackage{fontspec}
\setmainfont{Latin Modern Roman}
\usepackage{amsmath,amssymb,mathtools,unicode-math}
\setmathfont{Latin Modern Math}
\usepackage{enumitem,needspace,xurl,hyperref}
\setlength{\parindent}{0pt}
\setlength{\parskip}{5pt}
\setlength{\emergencystretch}{3em}
\newcommand{\N}{\mathbb{N}}
\newcommand{\Z}{\mathbb{Z}}
\newcommand{\Q}{\mathbb{Q}}
\newcommand{\R}{\mathbb{R}}
\newcommand{\C}{\mathbb{C}}
\newcommand{\F}{\mathbb{F}}
\newcommand{\A}{\mathbb{A}}
\newcommand{\PP}{\mathbb{P}}
\newcommand{\E}{\mathbb{E}}
\newcommand{\eps}{\varepsilon}
\newcommand{\dd}{\,\mathrm{d}}
\begin{document}
\section*{UnsolvedMath ID: 20001641}
\textbf{Source problem code:} catalogue entry
\subsection{Definitions and mathematical statement}
The canonical field is a section-boundary fragment, not an independent mathematical question. The aid identifies ``Problem 1.1?'' as the missing tail of Problem 1.2, followed by the heading \emph{Geodesics in Outer Space} and an introduction to Section 2. The text attributes to Juan Souto a discussion comparing Teichm\"uller and Weil--Petersson geodesic analogies with Lipschitz geodesics in Outer Space, and announces a multipart question that is not contained here. The actual Problem 2.1 begins in another record. No missing parts are invented.

The aid gives a distinct mathematical supplement about reverse costs. Let \(d\) be an asymmetric metric, let \(\gamma:I\to X\) satisfy \(d(\gamma(s),\gamma(t))=t-s\) for \(s<t\), and define
\[
d_{\max}(x,y)=\max\{d(x,y),d(y,x)\},\qquad
 d_{\Sigma}(x,y)=d(x,y)+d(y,x).
\]
A quasi-geodesic for one of these symmetrizations means constants \(K\ge1,C\ge0\) give the usual two-sided linear distance bounds in \(|t-s|\). The aid's equivalence is that this unit forward path is a quasi-geodesic for either symmetrization exactly when
\[
\exists A,B\ge0\ \forall s<t\text{ in }I,\qquad
 d(\gamma(t),\gamma(s))\le A(t-s)+B.
\]
The forward distance already gives the lower quasi-geodesic bound; the reverse-cost condition supplies the missing upper bound. Constants must be uniform over all pairs of times.

For volume-one marked Outer Space \(\mathrm{CV}_N\), \(N\ge2\), points are finite marked metric graphs with fundamental group \(F_N\). Let \(\ell_X(\alpha)\) be the length of a conjugacy class in \(X\), and let \(d_L(X,Y)\) be the logarithm of the optimal marking-preserving Lipschitz constant. The finite-candidate length-ratio formula gives, with \(X=\gamma(s)\), \(Y=\gamma(t)\),
\[
d_L(Y,X)=\log\max_{\alpha\in\operatorname{Cand}(Y)}\frac{\ell_X(\alpha)}{\ell_Y(\alpha)}.
\]
The finite source-graph candidate family consists of the standard simple-loop, figure-eight and barbell types. Thus reverse-cost control is equivalent to
\[
\ell_{\gamma(s)}(\alpha)\le e^{A(t-s)+B}\ell_{\gamma(t)}(\alpha)
\quad\text{for every }\alpha\in\operatorname{Cand}(\gamma(t)),\quad s<t.
\]
Candidates are taken at the later endpoint, which is the source of the reverse comparison, not automatically at the earlier endpoint.

The aid's sufficient budget criterion combines thick-part quasi-symmetry with control on thin pieces. On a fixed thick part, shortest nontrivial loop lengths have a uniform positive lower bound and known quasi-symmetry bounds reverse cost by a uniform multiple of forward duration. If each interval of the path can be subdivided so that the remaining thin pieces have aggregate reverse cost at most \(A(t-s)+B\), uniformly over intervals, the asymmetric triangle inequality gives a linear global reverse budget. An additive allowance on each separate thin piece without controlling their aggregate is not the same hypothesis.

The rank-two thin-rose warning is explicit in compatible coordinates: take the same marking on roses with loop lengths \((1/2,1/2)\) and \((\varepsilon,1-\varepsilon)\), \(0<\varepsilon<1/2\). Their forward distance from the balanced rose to the thin rose is \(\log(2(1-\varepsilon))<\log2\), while the reverse distance is \(\log(1/(2\varepsilon))\to\infty\). Bounded forward duration alone therefore cannot provide uniform reverse control. This supplement is not the missing multipart source problem.

The registry labels the entry open/open triage; attempt 1 is context only, with problem status not a problem, internal rigor proved and novelty confidence low. The aid is machine generated, not peer reviewed, and requires independent verification before citation.
\subsection{Short English statement}
The stored text is an introduction at a PDF section boundary, with no complete question. The aid's separate supplement identifies the extra condition needed for a forward Lipschitz geodesic to remain a quasi-geodesic after symmetrization: a uniform linear reverse-cost bound, equivalently exponential reverse candidate-length bounds. It does not reconstruct the absent multipart problem.
\subsection{Sources}
\begin{itemize}
\item Local UnsolvedMath registry, permanent ID \(20001641\), source code \texttt{catalogue entry}; \texttt{aim-geometric-group-theory-notes.json}, \texttt{notes[149]}.
\item AIM, \emph{The geometry of the outer automorphism group of a free group}, boundary of Problems 1.2 and Section 2, \emph{Geodesics in Outer Space}, discussion attributed to Juan Souto: \url{https://aimath.org/WWN/outerauto/outerauto.pdf}. The boundary identification is the aid's report, not independent PDF inspection.
\item Supplied research attempt 1 and registry update August 14, 2026. The original literature field is empty; no full authoritative references for the candidate formula or thick-part quasi-symmetry are supplied.
\end{itemize}
\end{document}
